\documentclass[../../../include/open-logic-section]{subfiles}

\begin{document}
\olfileid[fr]{sth}{card-arithmetic}{expotough}

\olsection[Quelques simplifications]{Quelques simplifications de l'exponentiation cardinale}

La définition de $\canonord$ était un peu compliquée,
mais elle nous a donné un résultat utile pour
l'addition et la multiplication cardinales,
\olref[simp]{cardplustimesmax}. L'exponentiation
transfinie, en revanche, ne se simplifie pas
aussi facilement. Pour expliquer pourquoi,
commençons par un résultat qui étend une règle
familière du cas fini, bien que sa démonstration
se situe à un niveau d'abstraction élevé :

\begin{prop}\ollabel{simplecardexpo}
$\cardexpo{\cardfont{a}}{\cardfont{b} \cardplus \cardfont{c}} =
\cardexpo{\cardfont{a}}{\cardfont{b}} \cardtimes
\cardexpo{\cardfont{a}}{\cardfont{c}}$ et
$\cardexpo{(\cardexpo{\cardfont{a}}{\cardfont{b}})}{\cardfont{c}} =
\cardexpo{\cardfont{a}}{\cardfont{b} \cardtimes \cardfont{c}}$,
pour tous cardinaux $\cardfont{a}, \cardfont{b},
\cardfont{c}$.
\end{prop}

\begin{proof}
Pour la première égalité, considérons une fonction
$f \colon
(\cardfont{b}\disjointsum\cardfont{c})
\to \cardfont{a}$. «~Découpons-la~» en définissant
$f_\cardfont{b}(\beta) = f(\beta, 0)$ pour
chaque $\beta \in \cardfont{b}$, et
$f_\cardfont{c}(\gamma) = f(\gamma, 1)$
pour chaque $\gamma \in \cardfont{c}$.
L'application
$f \mapsto \tuple{f_\cardfont{b},f_\cardfont{c}}$\footnote{
Note de l'édition française : l'original écrit ici
$f_\cardfont{b} \times f_\cardfont{c}$ ;
l'image doit être le couple des deux fonctions,
élément du produit de leurs espaces de fonctions.}
est !!a{bijection}
$\funfromto{\cardfont{b} \disjointsum \cardfont{c}}
{\cardfont{a}}
\to (\funfromto{\cardfont{b}}{\cardfont{a}}
\times \funfromto{\cardfont{c}}{\cardfont{a}})$.

Pour la seconde égalité, considérons une fonction
$f \colon \cardfont{c} \to
(\funfromto{\cardfont{b}}{\cardfont{a}})$ ;
pour chaque $\gamma \in \cardfont{c}$,
nous avons donc une fonction
$f(\gamma) \colon \cardfont{b} \to \cardfont{a}$.
Définissons alors
$f^*(\beta, \gamma) = (f(\gamma))(\beta)$
pour chaque
$\tuple{\beta, \gamma} \in
\cardfont{b} \times \cardfont{c}$.
L'application $f \mapsto f^*$ est
!!a{bijection}
$\funfromto{\cardfont{c}}
{(\funfromto{\cardfont{b}}{\cardfont{a}})}
\to \funfromto{\cardfont{b} \times \cardfont{c}}
{\cardfont{a}}$.\footnote{Note de l'édition
française : l'original met
$\cardfont{b}\cardtimes\cardfont{c}$
à la place de
$\cardfont{b}\times\cardfont{c}$
comme domaine de $f^*$.
La fonction construite agit sur les couples
du produit cartésien ; la cardinalité de ce
produit est
$\cardfont{b}\cardtimes\cardfont{c}$.}
\end{proof}

Nous aimerions maintenant disposer d'un moyen
simple de calculer
$\cardexpo{\cardfont{a}}{\cardfont{b}}$
lorsque les cardinaux sont infinis.
Voici un premier pas dans cette direction :

\begin{prop}\ollabel{cardexpo2reduct}
Si $2 \leq \cardfont{a} \leq \cardfont{b}$
et si $\cardfont{b}$ est infini, alors
$\cardexpo{\cardfont{a}}{\cardfont{b}} =
\cardexpo{2}{\cardfont{b}}$.
\end{prop}

\begin{proof}
\begin{align*}
	\cardexpo{2}{\cardfont{b}} &\leq
	\cardexpo{\cardfont{a}}{\cardfont{b}}
	\text{, puisque $2 \leq \cardfont{a}$}\\
	&\leq \cardexpo{(2^\cardfont{a})}{\cardfont{b}}
	\text{, par le \olref[opps]{lem:SizePowerset2Exp}}\\
	&= \cardexpo{2}{\cardfont{a} \cardtimes \cardfont{b}}
	\text{, par la \olref{simplecardexpo}}\\
	&= \cardexpo{2}{\cardfont{b}}
	\text{, par le \olref[simp]{cardplustimesmax}}
\end{align*}
\end{proof}

Il ne faut guère espérer simplifier davantage,
car $\cardfont{b} < \cardexpo{2}{\cardfont{b}}$
par \olref[card-arithmetic][opps]{cantorcor}.\footnote{
Note de l'édition française : l'original renvoie ici
au lemme qui identifie
$\card{\Pow{A}}$ et $2^{\card{A}}$ ;
l'inégalité stricte est formulée dans le corollaire
de Cantor cité ici.}
Cela ne nous dit cependant pas quoi conclure
pour $\cardexpo{\cardfont{a}}{\cardfont{b}}$
si $\cardfont{b} < \cardfont{a}$.
Bien sûr, si $\cardfont{b}$ est \emph{fini},
nous savons comment procéder.

\begin{prop}
Si $\cardfont{a}$ est infini et si
$0 < n \in \omega$, alors
$\cardexpo{\cardfont{a}}{n} = \cardfont{a}$.\footnote{
Note de l'édition française : l'original suppose seulement
$n \in \omega$ ; pour $n=0$, on a
$\cardexpo{\cardfont{a}}{0}=1$.
Il faut donc exclure ce cas.}
\end{prop}

\begin{proof}
$\cardexpo{\cardfont{a}}{n}
= \cardfont{a} \cardtimes \cardfont{a}
\cardtimes \ldots \cardtimes \cardfont{a}
= \cardfont{a}$,
par le \olref[simp]{cardplustimesmax}.
\end{proof}
\noindent
De plus, dans certains autres cas, nous pouvons
déterminer la taille de
$\cardexpo{\cardfont{a}}{\cardfont{b}}$ :

\begin{prop}
Si $2 \leq \cardfont{b} < \cardfont{a} \leq
\cardexpo{2}{\cardfont{b}}$ et si $\cardfont{b}$
est infini, alors
$\cardexpo{\cardfont{a}}{\cardfont{b}}
= \cardexpo{2}{\cardfont{b}}$.
\end{prop}

\begin{proof}
$\cardexpo{2}{\cardfont{b}}
\leq \cardexpo{\cardfont{a}}{\cardfont{b}}
\leq
\cardexpo{(\cardexpo{2}{\cardfont{b}})}{\cardfont{b}}
= \cardexpo{2}{\cardfont{b}\cardtimes\cardfont{b}}
= \cardexpo{2}{\cardfont{b}}$,
en raisonnant comme dans la \olref{cardexpo2reduct}.
\end{proof}
\noindent
Au-delà de ces cas, les choses deviennent
nettement plus subtiles.

\end{document}
